\documentclass{homework}
\author{Morales, Samantha}
\date{\today}
\class{CSCI 3203: Tashfeen's Machine Learning I}
\title{Homework 7}
\address{
  Computer Science,
  Petree College of Arts \& Sciences,
  Oklahoma City University
}
\begin{document} \maketitle

\question $(b^{\ell},\p b^{(\ell)} / \p \ep)$ must have the same dimensions. Show the following pairs have the same dimensions,
\[
	\left(b^{(2)},\frac{\p \ep}{\p b^{(2)}}\right), \quad \left(W^{2}, \frac{\p \ep}{\p W^{(2)}}\right), \quad \left(b^{(1)}, \frac{\p \ep}{\p b^{(1)}}\right), \quad \left(W^{(1)}, \frac{\p \ep}{\p W^{(1)}}\right)
\]

\question Use the chain rule to expand the following derivative,
\[
	\frac{df_n(f_{(n-1)}(f_{(n-1)}(...f(x))))}{dx}
\]

\question Use the follwoing recursive definition,
\[
	\frac{\p\ep}{\p b^{(\ell)}} = W^{(\ell + 1)^\dagger} \left(\frac{\p\ep}{\p b^{(\ell + 1)}}\right) \circ \sigma' (z^{(\ell)})
\]

\question Look up "Universal approximation theorem (UAT)" on the internet. Wikipedia is a good starting point. According to the UAT, multilayer feed-forward networks with as few as $n$ hidden later(s) are universal approximators. What is the value of $n$?

\begin{sol}
$n = 1$
\end{sol}

\question Explore the Anderson's Iris data set. Take 15 setosas, 15 versicolors and 15 virginicas for your testing set. One-hot encode the training labels and you might want to use np.expand\_dims(X, axis=-1) to turn the rows of a matrix X into column vectors. This is however dependant on your implementation. 
\begin{enumerate}[label=(\arabic*)]
	\item What is the baseline accuracy of the testing set?
	\item Draw a scatter plot of sepal width (cm) against the sepal length (cm). Colour each data point as setosa, versicolor, or virginica and label the legends appropriately.
	\item Write a neural network from scratch to classift the flowers. How did you pick the hyper parameters? 
	\item Give the final values of the following, 
	\begin{enumerate}
		\item Number of hidden layers and number of neural units in each hidden layer.
		\item The learning rate $n$.
		\item The stopping criteria for training.
		\item The strategy for initialising the weights and biases
		\item The training error curve.
		\item The final testing accuracy.
	\end{enumerate}
\end{enumerate}
\end{document}
